
\section*{摘要}
\addcontentsline{toc}{section}{摘要:}

本文围绕多机械臂搬运场景中的任务调度与运行模式选择问题展开研究。针对圆形作业区域内多类待搬运物块、不同类型收集盒以及双臂协同搬运需求，本文将实际机械臂搬运过程抽象为一个具有任务分配、执行排序、资源互斥和协同约束的运筹学优化问题。研究对象不是机械臂的底层运动控制或动力学建模，而是在给定任务、机械臂数量、作业时间和能耗参数的条件下，合理安排每个物块的执行机械臂、搬运顺序和协同模式，使系统整体完成时间较短、能耗较低，并兼顾机械臂之间的负载均衡。

在模型构建方面，本文以物块搬运任务为基本调度单元，考虑单臂搬运任务与双臂协同搬运任务的差异，建立了包含任务选择变量、开始时间变量、完成时间变量和模式选择变量的整数规划模型。目标函数以最小化最大完成时间为主，同时引入能耗指标和负载均衡指标用于评价不同调度方案的综合性能。针对该问题的组合复杂性，本文设计并实现了四类求解方法，包括 CP-SAT 求解器、融合剪枝的隐枚举法、启发式算法以及热启动后的隐枚举法，并通过多组算例对不同方法的求解效果和计算效率进行了比较。

在此基础上，本文进一步引入非线性优化分析，对机械臂速度倍率与时间、能耗之间的关系进行建模，研究不同速度选择对搬运效率和能耗代价的影响。同时，本文结合 KKT 条件对非线性优化结果进行可行性和最优性验证，并通过二臂与三臂场景的对比分析，讨论增加机械臂数量后系统在完成时间、能耗和模式选择方面的变化规律。实验结果表明，融合剪枝的隐枚举法能够在保持较高求解质量的同时减少无效搜索，启发式算法具有较好的计算速度，热启动后的隐枚举法能够进一步利用已有可行解提高搜索效率；非线性优化部分则说明机械臂速度调整会带来时间收益与能耗代价之间的权衡，为多机械臂调度方案的综合评价提供了补充依据。


\noindent {\heiti 关键词：}
多机械臂调度；运筹学优化；整数规划；隐枚举法；分支定界；启发式算法；非线性优化；KKT 条件